\section{Further questions}
\label{sec:questions}

We finish by discussing some natural follow-up questions and conjectures.

\subsection{Open surfaces}
A natural question is whether our results extend to open surfaces, i.e.\ compact surfaces with boundary or finite-type surfaces with punctures.

For such surfaces, geodesic currents admit at least two natural definitions, depending on how one treats the ends. 
If the ends are realized as cusps, one obtains a space $\Curr_{\mathrm{cusp}}(S)$, defined analogously to Definition~\ref{def:currents} as the space of $\pi_1(S)$-invariant 
Radon measures on the space of geodesics in the universal cover (see, e.g.,~\cite[Definition~2.2]{Sas22:Currents}). In this setting, the space of geodesics includes geodesic arcs connecting cusps. 
If instead one restricts to geodesics projecting to the convex core,
one obtains a space denoted $\Curr_{\mathrm{open}}(S)$ (\cite[Section~2.6]{DLR10:DegenerationFlatMetrics}). See
Figure~\ref{fig:open_surfaces}.

\begin{figure}
\centering
\begin{subfigure}{0.48\textwidth}
\centering
\resizebox{60mm}{!}{
\svgfigbig{cusp}}
\caption{Cusped surface with an arc $A$ and a closed geodesic $C_n$
  wrapping around the arc. We have $C_n \to 2A$.}
\end{subfigure}
\hfill
\begin{subfigure}{0.48\textwidth}
\centering
\resizebox{60mm}{!}{
\svgfigbig{open}}
\caption{Open surface with two boundary parallel curves $B_1$ and $B_2$. The sequence of curves $C_n$ also sits in this surface, and we have $\frac{1}{n}C_n \to B_1 + B_2$.}
\end{subfigure}
\caption{Examples of cusped and open surfaces, and currents on them.}
\label{fig:open_surfaces}
\end{figure}

In $\Curr_{\mathrm{cusp}}(S)$, the intersection pairing fails to be
continuous; see~\cite[Remark~2.26]{MGT21:Smoothings} and
also~\cite{Sas22:Currents, T22:ThurstonCompactification}. It is thus
hard to get started working entirely in this setting.


The case of $\Curr_{\mathrm{open}}(S)$ is different. Even though intersection number is continuous in this case, it is not true that our defining properties of $\AConvex(S)$ guarantee a dual geodesic current.
The key issue is that the intersection number is degenerate: any
geodesic current $\mu$ supported on boundary curves will have zero
intersection number with all geodesic currents $\nu \in \Curr(S)$.

To better understand the relation between the cusped and open
situation, consider Figure~\ref{fig:open_surfaces} in a surface with
two ends for simplicity.
In the cusped surface, let $A$ be the simple arc going from cusp to cusp, which is in $\Curr_{\mathrm{cusp}}(S)$.
Let $C_n$ be a sequence of closed geodesics obtained by following $A$
from near one cusp to the other, wrapping around the second cusp $n$ times,
going back along $A$, and wrapping around the first cusp $n$ times
before closing up. In the weak$^*$-topology of
$\Curr_{\mathrm{cusp}}(S)$, we have $\lim_{n\to\infty}C_n = 2A$. (See, for example,
Sasaki~\cite[Section~6]{Sas22:Currents}.)

In the $\Curr_{\mathrm{open}}$ case, let $B_1$ and $B_2$ be the
boundary parallel curves homotopic to each respective boundary
component. The sequence of curves $C_n$ can also be considered in
$\Curr_{\mathrm{open}}(S)$. Without scaling, $C_n$ does not converge
to any current, but it does projectively. In fact, $\lim_{n \to
  \infty} \frac{1}{n} C_n = B_1 + B_2$.

Consider the geometric intersection number between arcs and curves in the cusped surface and curves in the open surface, i.e., $i \colon \Curves_{\mathrm{cusp}}(S) \times \Curves_{\mathrm{open}}(S)$.
We claim
\[
i(C_n, \tfrac{1}{n} C_n) \to i(2A, B_1 + B_2).
\]
Indeed, $i(C_n, C_n) = 4n$, and $i(2A, B_1+B_2)=4$, from which the result follows.
Following this example, we conjecture there is a well-defined and
continuous geometric intersection number between
$\Curr_{\mathrm{open}}(S)$ and $\Curr_{\mathrm{cusp}}(S)$.

\begin{conjecture}
There exist a continuous function 
\[
I \colon \Curr_{\mathrm{cusp}}(S) \times \Curr_{\mathrm{open}}(S) \to \mathbb{R}
\]
extending the natural geometric intersection pairing between $\Curves_{\mathrm{cusp}}(S)$ and $\Curves_{\mathrm{open}}(S)$.
\end{conjecture}

It is known that if one restricts the right-hand-side
factor to currents supported on a preimage of a compact subset on the
surface, the intersection number is continuous
(See \cite[Theorem~2.4]{BIPP21:Systoles}
and~\cite[Section~6]{Sas22:Currents}).

We can use this intersection number to propose a duality. Say that a
function $f \colon \Curves_{\mathrm{cusp}}(S) \to \mathbb{R}$
satisfies \emph{smoothing} if it satisfies smoothing between any
closed multi-curves, as well as smoothing between arcs and curves and
between arcs, in the sense that after resolving any essential crossing
(any intersection that cannot be homotoped away), the function does
not increase.

Similarly, say that a function
$f \colon \Curves_{\mathrm{open}}(S) \to \mathbb{R}$ satisfies
smoothing if it does not increase after resolving any essential
crossing.

For examples of such functionals, consider an embedded graph $\mathcal{G} \subset S'$.
Equipping $\mathcal{G}$ with a metric, we get a notion of length of
$\ell_\Sigma$ curves in $\Curves_{\mathrm{open}}(S)$ with respect to
this embedded metric graph. (See also
Section~\ref{subsec:embedded_graphs}.) When $\mathcal{G}$ is a rose-graph
(embedded and
with one vertex), Erlandsson~\cite{Erlandsson:WordLength} has
constructed examples of geodesic currents $\mu_{\mathcal{G}}$ in
$\Curr_{\mathrm{open}}(S)$ (in fact, half-integer weighted arcs and
multi-curves) so that
\[
\ell_{\Sigma}(C)=I(\mu_{\mathcal{G}}, C).
\]
This motivates the following question.

\begin{question}\label{quest:open-to-cusp}
Given a function $f \colon \Curves_{\mathrm{open}}(S) \to \mathbb{R}$
which is additive and satisfies smoothing, are there additional
assumption guaranteeing the existence of a unique geodesic current $\mu \in f \colon \Curr_{\mathrm{cusp}}(S)$ so that
\[
f(C)=i(C, \mu)
\]
for every $C \in \Curves_{\mathrm{open}}$?
\end{question}

Unfortunately, additivity and smoothing alone are not enough. For
example, take the function
$f \colon \Curves_{\mathrm{open}}(S) \to \mathbb{R}$ assigning
hyperbolic length to non-boundary parallel curves and $0$ to
boundary-parallel curves. This functional still satisfies smoothing
(since no other curve can cross a boundary-parallel curve essentially)
but does not have a dual geodesic current in
$\Curr_{\mathrm{cusp}}(S)$, since by continuity of $I$ it would have
to assign positive length to limits of curves wrapping around the
boundary. See~\cite[Remark~2.26]{MGT21:Smoothings} for details.

In the other direction, we have a parallel conjecture.

\begin{conjecture}\label{conj:cusp-to-open}
Given a function $f \colon \Curves_{\mathrm{cusp}}(S) \to \mathbb{R}$ which is additive and satisfies smoothing, there exists a unique geodesic current $\mu \in f \colon \Curr_{\mathrm{open}}(S)$ so that
\[
f(C)=i(C, \mu)
\]
for every $C \in \Curves_{\mathrm{cusp}}$(S).
\end{conjecture}

We have two motivating examples of functionals for
Conjecture~\ref{conj:cusp-to-open}.

The first one is related to the example on graphs above. Take again an
embedded metric graph $\mathcal{G} \subset S$, and define the
intersection $i(\mathcal{G},C)$ to be the minimal transversal
intersection between homotopy classes of embedded representatives of
$\mathcal{G}$ into $S$ and a homotopic representative of $C$, for
every closed curve in $\Curves_{\mathrm{cusp}}(S)$. This defines a
function $f \colon \Curves_{\mathrm{cusp}}(S) \to \mathbb{R}$. Then
our conjecture would assert the existence of a geodesic current
$\mu \in \Curr_{\mathrm{open}}(S)$. Presumably, as in Erlandsson's
result, $\mu$ in this case consists of half-integral combination of
curves.

As a second example, consider a flat metric on the completion of $S'$ so that the missing points become conical singular points. Now the length of an arc between cusps is the length of the arc from cone point to cone point, which is a finite number.
This gives a well-defined notion of length $\ell \colon \Curves_{\mathrm{cusp}}(S) \to \mathbb{R}$.
On the other hand,  $\Curr_{\mathrm{open}}(S)$ are geodesic currents
whose support is contained in the closure of the subspace of geodesics
that do not go through cone points. The natural Liouville metric for
singular flat metrics is an example of such a
current~\cite[Section~3.2]{BL17:RigidityFlat},
and seems to be dual
to~$\ell$.


\subsection{Generalizations to hyperbolic groups}
In recent work of Cantrell and Reyes and, independently, Kapovich and the first author, the following property has played a key role in extending stable lengths coming from group actions of hyperbolic groups $\Gamma$ acting on uniformly quasi-geodesic hyperbolic spaces $X$~\cite{KMG:BBP,CR25:Manhattan, CMGR26:GreenMetrics}.
Roughly, a Gromov hyperbolic space is a metric space whose triangles are thin. A group $\Gamma$ is said to be Gromov hyperbolic if its Cayley graph is.

A metric space $X$ is uniformly quasi-geodesic if there exists
constants $\lambda >0$ and $\epsilon \geq 0$, so that
every two points $x, y \in X$ are connected by a
$(\lambda,\epsilon)$-quasi-geodesic. We say that an action
$\Gamma \acts X$ of a non-elementary hyperbolic group~$\Gamma$ on a uniformly
quasi-geodesic hyperbolic metric space~$(X,d)$ has the \emph{bounded
  backtracking property} (see~\cite[Section~6.1]{CR25:Manhattan}) if given an orbit map
$\iota \colon \Cay(\Gamma) \to X$, there exists a constant $C>0$ so that
for every two points $x,y \in \Cay(\Gamma)$, and geodesic $[x,y]$,
there exists some quasi-geodesic $[\iota(x),\iota(y)]$ in $X$ so that
$\iota([x,y])$ and $[\iota(x),\iota(y)])$ are $C$-Hausdorff close in $d$.

There are examples of stable lengths satisfying bounded backtracking
which do not satisfy quasi-smoothing
(see~\cite[Section~9]{CMGR26:GreenMetrics}). On the other hand, one
can ask if all functions on $\pi_1(S)$ satisfying quasi-smoothing
arise as the stable length of some action of $\pi_1(S)$ satisfying
bounded backtracking. 

\begin{question}
Let $\Gamma$ be a surface group, and let $f \colon \Gamma \to \mathbb{R}$ be a symmetric and stable functional satisfying quasi-smoothing (and hence, extending continuously to geodesic currents).
Does there exists a uniformly quasi-geodesic Gromov hyperbolic space $X$ on which $\Gamma$ acts with the bounded backtracking property and so that $\ell_X = f$, where $\ell_X$ denotes the stable length of the action $\Gamma \acts X$?
\end{question}

Motivated by this, it would be interesting to generalize the quasi-smoothing and smoothing condition to the general setting of hyperbolic groups.

\begin{question}
What is a generalization of the properties of smoothing and
quasi-smoothing for functionals on non-elementary hyperbolic groups
that makes the answer to the
above question work in that setting?
\end{question}

Related to Theorems~\ref{thm:intersections_group}
and~\ref{mainthm:skora}, it would be interesting to give a list of
axioms, in the spirit of oriented smoothing, characterizing
small actions of a non-elementary hyperbolic group $\Gamma$.

\begin{question}
\label{qu:generalize_intersections}
Let $\Gamma$ be a non-elementary Gromov hyperbolic group action on an
$\mathbb{R}$-tree $T$. Give axioms  characterizing when the action is small, i.e., has
virtually cyclic edge stabilizers.
\end{question}

\subsection{Quasi-morphisms} 
\label{sub:quasimorphism} 

The hypotheses of \emph{oriented quasi-smoothing} are reminiscent of, but weaker than, the conditions for a map 
$f \colon \pi_1(S) \to \mathbb{R}$ to be a \emph{quasi-morphism}, that is, there exists $D > 0$ such that 
\[
|f(ab) - f(a) - f(b)| < D \quad \text{for all } a,b \in \pi_1(S).
\]
A quasi-morphism is called \emph{homogeneous} if $f(a^n) = n f(a)$ for every $a \in \pi_1(S)$ and $n \ge 1$.

Quasi-morphisms have played a prominent role in the study of stable commutator length and bounded cohomology of groups~\cite{Cal09:scl,Pic97:BoundedCohomology_currents}.

It would be interesting to understand whether one can obtain homogeneous quasi-morphisms from curve functionals satisfying oriented quasi-smoothing and stability.

\begin{question}
Let $f \colon \pi_1(S) \to \mathbb{R}$ satisfy oriented quasi-smoothing and stability.  
Are there additional conditions on $f$ ensuring that the function
\[
g(a) \coloneqq f(a) - f(a^{-1})
\]
is a homogeneous quasi-morphism?
\end{question} 

To get a non-trivial function~$g$, the functional~$f$ must be
asymmetric. See Appendix~\ref{sec:asymmetric} for examples.

% \subsection{Functionals valued on ordered abelian groups}

% We expect most of the results in this paper to extend, with only minor modifications, to the setting of $\Lambda$-valued functionals and $\Lambda$-valued geodesic currents, where $(\Lambda, <)$ denotes an ordered abelian group---that is, an abelian group equipped with a total order satisfying $a < b$ and $c \le d$ imply $a + c < b + d$.
% This generalization is natural in the context of the theory of $\Lambda$-trees~\cite{Sko90:GeometricAction} and has already proved fruitful in the study of compactifications of character varieties of surface-group representations into higher-rank Lie groups~\cite{BIPP24:PositiveCrossratios}.
% We plan to explore these extensions and their applications in future work.

\subsection{Green metrics}
\label{subsec:green}
Let $\Gamma=\pi_1(\Sigma) \subset \Isom^+(\mathbb{H})=\PSL_2(\mathbb{R})$ be a countable, discrete cocompact subgroup, for some hyperbolic metric $\Sigma$ on $S$.
Let $\mu$ be a finitely supported probability measure on $\Gamma$ whose support contains a symmetric generating set of $\Gamma$.
Let $w_n \coloneqq g_1 \cdots g_n$ be the random walk with respect to $\mu$, where each $(g_i)$ is identically and independently distributed with distribution $\mu$. Let $(\Gamma^{\mathbb{N}}, \mu^{\mathbb{N}})$ be the step space, and the map $\pi \colon \Gamma^{\mathbb{N}} \to \Gamma^{\mathbb{N}}$ as $\pi((g_n)_{n \in \mathbb{N}}) \coloneqq (w_n)_{n \in \mathbb{N}}$. The image of $\pi$ is equipped with the measure $(\mathbb{P}_{\mu})_*(\mu^{\mathbb{N}})$.
Define the metric \[ d_{\mu}(x,y) \coloneqq - \log \mathbb{P}_{\mu}(\exists n : w_n x =y).\]
In words, we are considering the probability that a random walk departing from $x$ will reach $y$ after some number of steps.
This metric has an associated stable length function $\ell_{\mu}(g) \coloneqq \lim_n d_{\mu}(x, g^n x)/n$, for any hyperbolic element $g \in \Gamma$, which can then be thought as a curve functional $\ell_{\mu} \colon \Curves^+(S) \to \mathbb{R}$. The following question emerged from conversations of the first author with Eduardo Reyes.

\begin{question}
Does $\ell_{\mu}$ satisfy the hypotheses of
Theorem~\ref{thm:intersect} for some admissible probability measure $\mu$ (and thus come from a geodesic current)?
\label{q:green-smoothing}
\end{question}

\begin{question}
Does $\ell_\mu$ further satisfy the
hyperbolic metric hypotheses in Theorem~\ref{thm:hyp-lam-char}(H) for some admissible
probability measure~$\mu$ (and thus come
from a hyperbolic metric)?
\label{q:Green}
\end{question}

If answers to either of these questions is negative for every $\mu$,
they would imply a negative answer to the following well-known question.

\begin{question}
Given $\Sigma$, is there a finitely-supported probability measure
$\mu$ on $\Gamma=\pi_1(\Sigma)$ so that the induced translation length
$\ell_{\mu}$ is dual to the Liouville current $\mathcal{L}_{\Sigma}$?
\label{q:singularity}
\end{question}

By work of Cantrell-Tanaka~\cite[Theorem~1.2]{CT25} and
Gouëzel-Mathéus-Maucourant \cite[Theorem~1.2]{GMM18:Entropy}, a
negative answer to this question would settle the so-called
\emph{singularity conjecture} for $\SL_2(\mathbb{R})$
(see~\cite[Conjecture, page~~259]{KLP11:Singularity}). In short, by
these results Question~\ref{q:singularity} is equivalent to asking if
there is a probability measure $\mu$ on $\Gamma$ so that the induced
`hitting measure' $\nu_{\mu}$ on $\partial \Gamma$ is absolutely
continuous with respect to the Lebesgue measure (coming from the
hyperbolic metric $\Sigma$) on $\partial \Gamma$.
Kosenko-Tiozzo~\cite{KT22:Inequality} have studied the metric $d_{\mu}$ in some cases of Fuchsian groups $\Gamma$ with centrally symmetric fundamental domains.
Even a positive answer to Question~\ref{q:Green} would be interesting, since then
Theorem~\ref{thm:intersect} would give a dual geodesic current for the
Green metric, and then one could study what types of geodesic currents
can arise.

\subsection{Smoothing vs multi-smoothing}
\label{subsec:multismoothing}
The hypotheses of smoothing in Equation~\eqref{eq:smoothing} of Theorem~\ref{thm:intersect} apply to one essential crossing at a time. There is a stronger condition than smoothing, where multiple smoothings at essential crossings are performed at once, and the function is required to be non-increasing under such operation.
We will refer to this condition as \emph{multi-smoothing}.
It is clear multi-smoothing implies smoothing.
It is easily checked that intersection numbers satisfy the multi-smoothing property (see, for example, \cite{AH21:EffectiveCounting3}).
What is not so clear is that functionals of $\AConvex(S)$ (that is, satisfying stability, homogeneity, smoothing and additivity) must satisfy multi-smoothing. This is, of course, a consequence Theorem~\ref{thm:intersect}.
Other curve functionals that are not in $\AConvex(S)$ also satisfy
multi-smoothing. For example, Arana-Herrera also shows
that extremal length satisfies
multi-smoothing \cite{AH21:EffectiveCounting3}.
This prompts the following question.

\begin{question}\label{quest:multi-smoothing}
If a curve functional satisfies homogeneity, stability, convex union,
and smoothing, does it necessarily satisfy multi-smoothing?
\end{question}

In this context, Arenas and Neumann-Coto
recently studied the related concept of \emph{taut smoothing}  \cite{ANC24:TautSmoothings}.

\subsection{Domination}
\label{subsec:domination}
In another direction from Section~\ref{subsec:multismoothing},
we say that a curve $C$ \emph{dominates} another curve $C'$ if $i(C,D) \geq
i(C',D)$ for every curve $D$. We say a functional $f$ satisfies
\emph{domination} if for
every pair of curves $(C,C')$ with $C$ dominating $C'$, we have
$f(C) \geq f(C')$.
It is easy to see that if $C'$ is obtained from $C$ by
multi-smoothing, then $C$ dominates $C'$. There are however other
pairs $(C,C')$ with $C$ dominating $C'$ but $C'$ not obtainable from
$C$ via multi-smoothing
(see~\cite[Example]{NC01:ShortestGeodesics}).
It is clear, by definition, that intersection numbers satisfy domination.
All in all, we have the following implications on a functional~$f$:
\[
  \text{Domination} \Rightarrow \text{Multi-smoothing} \Rightarrow \text{Smoothing}.
\]

From Theorem~\ref{thm:intersect} we obtain the following weaker
theorem which is nevertheless conceptually interesting.

\begin{theorem}
A curve functional satisfies domination and additivity if and only if it is dual to a geodesic current.
\end{theorem}

This can be interpreted as saying that linear maps that preserve the poset structure of the space of multi-curves correspond to intersection numbers.

For similar reasons as for the smoothing (or multi-smoothing)
condition, the additivity hypothesis is necessary, since
the square
root extremal length can be easily seen to satisfy domination,
homogeneity and stability, but not additivity, and it is not dual to
an intersection number.
This suggests the following strengthening of
Question~\ref{quest:multi-smoothing}.

\begin{question}
If a curve functional satisfies homogeneity, stability,
convex union, and smoothing, does it necessarily satisfy domination?
\end{question}

\subsection{Convexity on train-track charts and local compactness}
\label{sec:convexity_charts}

One of the reasons multi-smoothing is a useful property is that,
coupled with convex union, homogeneity, and stability, it
yields convex functions when restricted to train-track coordinates.
Precisely, let $f$ be a curve functional satisfying the aforementioned
properties, including multi-smoothing. Let $\tau$ be a maximal train-track (i.e., a
train-track all of whose complementary components are trigons), let
$V(\tau)$ be the set of admissible weights on~$\tau$, $\ML(\tau)$ be the subspace of
measured laminations carried by $\tau$, and
$\phi_{\tau} \colon V(\tau) \to \ML(\tau)$ be the induced
chart on measured laminations.
Since $f$ extends continuously to geodesic currents \cite[Theorem~A]{MGT21:Smoothings}, we can consider the composition $f_{\tau} \coloneqq f \circ \phi_{\tau} \colon V(\tau) \to \mathbb{R}$. Under these assumptions, the same proof as in~~\cite[Appendix~A]{Mir04:SimpleGeodesics}, \cite[Theorem~2]{Thu16:Research}, \cite[Theorem~4.7]{AH21:EffectiveCounting3}, shows $f_{\tau}$ is convex and homogeneous. 

These curve functionals form a subset of the functionals that extend
continuously to currents, by~\cite[Theorem~A]{MGT21:Smoothings}.
Moreover, the map $\Psi \colon \Curr(S) \to \mathbb{R}^{\Curves(S)}$
defined by $\Psi(\mu) = i(\mu, \cdot)$ is a homeomorphism onto its
image when $\mathbb{R}^{\Curves(S)}$ is equipped with the product
topology,
by~\cite[Theorem~11]{DLR10:DegenerationFlatMetrics}.
Thus, a restatement of Theorem~\ref{thm:intersect} is as follows.

\begin{theorem}
   The image of~$\Psi$ is homeomorphically $\AConvex(S) \subset \bbR^{\Curves(S)}$.
\end{theorem}

It is a consequence of this last result and the local compactness of
the space of geodesic currents (recalled as
Proposition~\ref{prop:local_compactness} below) that the space
$\AConvex(S)$ with the product topology is locally compact, which is
not obvious \textit{a priori} by looking directly at the functionals.

The sublevel set $\Lambda_f \coloneqq \{ \vec{x} \in V(\tau)
: f_{\tau}(\vec{x}) \leq 1 \}$ is a compact convex subset of
$\mathbb{R}^N$, for some $N>0$ depending on the combinatorics of~$\tau$.
For any sequence of geodesic currents $\mu_i$ limiting to~$\mu$,
let $f_i(\cdot) = i(\mu_i, \cdot)$ for $f(\cdot) = i(\mu, \cdot)$.
Then $\Lambda_{f_i} \to \Lambda_f$ in the Hausdorff topology. With
respect to this topology,
the subset of compact convex subsets of $\mathbb{R}^N$ is locally
compact \cite[Chapter~6]{Lay92:Convex}. We see this phenomenon as a
shadow of the local compactness of $\AConvex(S)$, which suggests that
the same local compactness for $\Convex(S)$ should hold.

\begin{conjecture}\label{conj:convex-compact}
The set $\Convex(S)$ with the product topology is locally compact.
\end{conjecture}

Some examples of interest for this conjecture would be extremal
lengths associated to conformal structures or to elastic graphs
\cite[Sections 4.8 and 4.9]{MGT21:Smoothings}.






\subsection{Analogy to Poincar\'e duality}
\label{subsec:poincare_duality}
Theorem~\ref{thm:intersect} gives a new definition of geodesic currents in terms of functionals on curves, by giving a bijection between geodesic currents and a subset of the space of curve functionals.
As mentioned in the introduction, one can view this theorem in analogy to Poincaré duality for closed surfaces.  Given a 1-cohomology class $\alpha^*$ of an $n$-dimensional manifold (represented by homomorphisms from $\pi_1(S)$ to $\mathbb{R}$), Poincaré duality associates an $(n-1)$-dimensional homology class $\alpha$ (represented by real linear combinations of $(n-1)$-cycles). Since $n=2$, these cycles can be expressed as linear combinations of simple closed curves.  

In our setting, the role of a 1-cohomology class is played by a curve
functional that satisfies smoothing, additivity, and
stability. Although it is not a homomorphism from
$\pi_1(S)$ to anything, it preserves a certain combinatorial structure
(see Section~\ref{subsec:domination}). Likewise, the role of the homology
class is taken by a geodesic current.

Our theorem can be thought of in analogy to $\alpha \cap \beta =
  \alpha^*(\beta)$, where $\cap$ denotes the cap product in homology (algebraic intersection), and thus relating cap product to the homology-cohomology pairing.

  This aligns with our theorem, which establishes a duality via the \emph{geometric} intersection number $i(\mu_f, C) = f(C)$, suggesting that the geometric intersection number plays the role of the cap product. 

  
In Poincaré duality, one also has the standard relation $\alpha^* \cup \beta^* = \alpha \cap \beta$, where $\cup$ denotes the cup product in cohomology.  
A natural question that arises then is the following.

\begin{question} What is the appropriate analog of the cup product?
  More precisely, can one define the intersection $i(\cdot,\cdot)$ of
  two geodesic
  currents directly via their dual curve functionals, and prove its
  continuity?
\end{question}

\subsection{Asymmetric analogues}
\label{sec:asymmetric-analogues}

In this paper, we characterized additive, unoriented curve functionals that
satisfy smoothing as exactly those coming from geodesic currents.
We could instead look at functionals on \emph{oriented} curves.

\begin{definition}
  Let $\AConvex^+(S)$ be the space of functionals on oriented curves
  $\Curves^+(S)$ that satisfy additivity, oriented smoothing
  (connected and disconnected), and stability, endowed with the
  product topology. Similarly, $\AConvex^+_{\ge 0}(S)$ is the subspace
  of positive functionals.
\end{definition}

Appendix~\ref{sec:asymmetric} can be summarized as saying that
$\AConvex^+(S)$ is not very related to the space $\Curr^+(S)$ of
oriented geodesic currents: the natural map $\mu \mapsto
\vec\imath(\mu,\cdot)$ is neither injective nor surjective, in what
appear to be large ways.

Nevertheless the space $\AConvex^+_{\ge 0}(S)$ may be of independent
interest. There are many examples of such functionals. For instance,
you can take an asymmetric metric on $S$ and obtain a functional
on curves as in Section~\ref{subsec:ineqintersection}. As in the
symmetric case, if there are rectifiable curves in each homotopy class
the result is automatically positive and satisfies oriented smoothing.

For a more concrete example, you can take an asymmetric generating set
for $\pi_1(S)$ as a monoid whose (directed) graph is simple in~$S$, as
in the example in Appendix~\ref{sec:asymmetric}.

\begin{question}\label{quest:AConvex+-compact}
  Is the space $\AConvex^+_{\ge 0}(S)$ a reasonable space? For
  instance, is it locally compact?
\end{question}

\begin{question}
  Does every element of $\AConvex^+_{\ge 0}(S)$ come from an
  asymmetric pseudo-metric in the style of
  Corollary~\ref{maincor:length_dual}?
\end{question}

In this oriented context, the restriction that the functional is
positive is non-trivial (unlike for unoriented functionals): for any
functional $f$ and element $\omega \in H^1(S; \bbZ)$, for a sufficiently
large constant $\lambda>0$ the combination $f + \lambda\omega$ will take negative values.

We could also relax the condition of smoothing to quasi-smoothing, and
ask similar questions. If we drop positivity and also consider
quasi-smoothings, then the space of functionals includes the much
larger space of quasi-morphisms (Section~\ref{sub:quasimorphism}), which an infinite dimensional Banach space and hence not locally compact (see, e.g.,~\cite{CRH26:Quasimorphisms}).

Recall that $\QConvex_K^+(S) \subset \QConvex^+(S)$ denotes the subspace of curve functionals on oriented multi-curves satisfying oriented quasi-smoothing with constant~$K$, homogeneity, stability and convex union.
$\QConvex_K^+(S)$ contains, for example, the space of homogeneous $K$-quasi-morphisms. The space of $K$-quasi-morphisms is easily seen to be locally compact by an inequality
\[
f(C) \leq K' \ell_F(C)
\]
where $\ell_F$ is the word-length induced by any finite generating set, and $K'$ is a constant depending only on a finite subset of curves.
The hypotheses of oriented $K$-quasi-smoothing are much weaker than those of $K$-quasi-morphisms.

\begin{question}\label{quest:QConvex+-compact}
Is $\QConvex_K^+(S)$ locally compact?
\end{question}

More generally, one could ask about local compactness for various weakenings of the the hypotheses defining the subset $\AConvex(S)$; Conjecture~\ref{conj:convex-compact} and Questions~\ref{quest:AConvex+-compact} and \ref{quest:QConvex+-compact} are some examples.

%%% Local Variables:
%%% mode: latex
%%% TeX-master: "Intersections"
%%% End:

